
% Autoría principal: Mario Zambrano.
% Protocolo de las mediciones realizadas el 20-11-2026.

\subsection{Entorno de ejecución}

El experimento se realizó el 20 de noviembre de 2026 en una computadora con Windows 11 Home
(versión 10.0.26200), procesador AMD Ryzen 7 7730U with Radeon Graphics y 15,34 GiB de
memoria RAM. Se utilizó Python 3.11.9 dentro de un entorno virtual para mantener las
dependencias del proyecto aisladas de otras instalaciones.

Las bibliotecas utilizadas fueron
\texttt{grpcio} y \texttt{grpcio-tools} 1.85.1,
\texttt{protobuf} 7.36.2, \texttt{pandas} 3.1.1
y \texttt{matplotlib} 3.11.2. Las versiones se
comprobaron antes de ejecutar las mediciones y
corresponden a las definidas en
\texttt{requirements.txt}.

La práctica contempla una arquitectura de tres
nodos independientes, denominados
\texttt{nodo-A}, \texttt{nodo-B} y
\texttt{nodo-C}, asociados a los puertos
5000, 5001 y 5002. Para la comparación de
rendimiento se utilizó el cliente identificado
como \texttt{nodo-A}, que estableció una
conexión con el servidor \texttt{nodo-B}
a través de \texttt{127.0.0.1:5001}.

Los mecanismos TCP y gRPC se evaluaron por separado, deteniendo el servidor anterior
antes de iniciar el siguiente. En ambos casos se redujo la salida por consola y
se desactivó el registro de eventos del servidor, mientras que los clientes
mantuvieron sus registros automáticos.

\subsection{Variables, condiciones y control experimental}

Se consideraron dos variables independientes: el mecanismo de comunicación y el tamaño de
la carga útil. Para el primer factor se evaluaron sockets TCP y gRPC con Protocol
Buffers; para el segundo se utilizaron mensajes de 100 y 10240 bytes.

Las cargas útiles se generaron mediante la repetición de un carácter ASCII. De esta
manera, el contenido enviado tuvo la longitud establecida para ambas
implementaciones. Los tamaños indicados corresponden únicamente al campo
\texttt{message} y no incluyen los bytes adicionales de JSON, Protocol Buffers ni
los encabezados de transporte.

La variable de respuesta principal fue la latencia de ida y vuelta (RTT), expresada
en milisegundos. Se midió el tiempo transcurrido durante la operación de
envío de una solicitud y recepción de su acuse, utilizando
\texttt{time.perf\_counter()}.

\begin{equation}
RTT_{\mathrm{ms}} =
(t_{\mathrm{fin}}-t_{\mathrm{inicio}})
\times 1000
\end{equation}

Como segunda métrica se calculó el rendimiento en mensajes por segundo. Para
cada corrida se dividió el número de mensajes confirmados entre la duración
total de dicha corrida:

\begin{equation}
R = \frac{N}{T}
\end{equation}

En esta expresión, $N$ representa la cantidad de mensajes confirmados y $T$
la duración de la corrida en segundos.

El benchmark reutilizó una conexión por mecanismo y tamaño de carga útil. Los
mensajes se enviaron secuencialmente: cada solicitud debía recibir su acuse
antes de continuar con la siguiente. Se estableció un tiempo de espera de
dos segundos por operación.

\subsection{Número de repeticiones y calentamiento}

La combinación de los dos mecanismos con los dos tamaños de mensaje produjo
cuatro condiciones experimentales.

Antes de cada condición se ejecutó una corrida de calentamiento de 100 envíos.
Estas observaciones quedaron excluidas de los archivos consolidados utilizados
para el análisis estadístico.

Después del calentamiento se realizaron cinco corridas medidas de 100 envíos,
equivalentes a 500 observaciones por condición.

\begin{table}[htbp]
\centering
\small
\begin{tabular}{lrrr}
\toprule
Mecanismo & Carga útil & Corridas & RTT medidos \\
\midrule
TCP  & 100 B   & 5 & 500 \\
TCP  & 10240 B & 5 & 500 \\
gRPC & 100 B   & 5 & 500 \\
gRPC & 10240 B & 5 & 500 \\
\midrule
\textbf{Total} & & \textbf{20} & \textbf{2000} \\
\bottomrule
\end{tabular}
\caption{Distribución de las mediciones del
experimento. Se excluyeron 400 envíos de
calentamiento.}
\label{tab:diseno_experimental}
\end{table}

Las corridas medidas finalizaron con 2000 acuses correctos. La revisión
posterior de los registros confirmó que cada condición contenía 500 RTT
positivos y válidos, distribuidos en cinco corridas de 100 mensajes.

\subsection{Prueba simultánea de tres clientes}

Además del benchmark secuencial, se realizaron pruebas funcionales utilizando tres nodos en localhost. Para ejecutarlas se utilizó \texttt{analysis/run\_lamport\_experiment.py}, que inicia tres procesos independientes. Cada proceso dispone de un servidor, un cliente y su propio reloj lógico de Lamport.

La comunicación se organizó en un anillo: el nodo-A envía solicitudes al nodo-B, el nodo-B al nodo-C y el nodo-C al nodo-A. Cada solicitud recibe un acuse de respuesta.

Como comprobación funcional se documentaron veinte intercambios por mecanismo, distribuidos entre los tres clientes en cantidades de 7, 7 y 6. Los registros correspondientes se conservaron en \texttt{data/lamport20/}, junto con los archivos de verificación y sus hashes SHA-256.

También se realizaron cuatro pruebas adicionales de concurrencia, combinando TCP y gRPC con cargas útiles de 100 y 10240 bytes. En cada condición se enviaron 300 solicitudes, distribuidas en partes iguales entre los tres clientes.

El lanzador utiliza barreras para coordinar el inicio de los servidores, el registro de los primeros envíos y la finalización. Este procedimiento permite estudiar eventos concurrentes, aunque no representa tráfico espontáneo ni demuestra que los envíos ocurrieran exactamente al mismo instante.

Las pruebas se ejecutaron en un entorno Linux diferente del equipo Windows utilizado para medir las latencias. Las cuatro capturas de 300 solicitudes fueron ejecuciones automatizadas asistidas, cuya procedencia y comandos se documentan en \texttt{data/concurrencia/PROCEDENCIA.md}. Las corridas personales de veinte intercambios se documentaron por separado mediante capturas.

Estos registros se utilizaron para verificar intercambios, acuses y relaciones causales mediante los relojes de Lamport. No se incorporaron a las 2000 mediciones RTT del benchmark, porque corresponden a condiciones y objetivos experimentales diferentes.


\subsection{Trazabilidad y protección de los datos}

Los resultados del benchmark se almacenaron en cuatro archivos CSV
dentro del directorio \texttt{data/}:

\begin{itemize}
    \item \texttt{latency\_sockets.csv}
    \item \texttt{runs\_sockets.csv}
    \item \texttt{latency\_grpc.csv}
    \item \texttt{runs\_grpc.csv}
\end{itemize}

Los archivos de latencia incluyen el mecanismo, tamaño del mensaje, corrida, intento, RTT, estado del acuse y fecha de registro. Los archivos de rendimiento contienen la duración de cada corrida, el número de mensajes confirmados y la tasa calculada en mensajes por segundo.

El análisis se realizó mediante \texttt{analysis/analisis\_resultados.py}, utilizando pandas y matplotlib. El programa comprobó la estructura de los datos, la cantidad de registros por condición y la ausencia de acuses fallidos o latencias inválidas. A partir de los archivos originales generó los resúmenes estadísticos, las tablas LaTeX y las figuras.

Los CSV consolidados se conservaron sin modificaciones manuales. También se calcularon sus huellas SHA-256, registradas en \texttt{data/SHA256SUMS}. La integridad se comprobó recalculando cada hash y comparándolo con su valor almacenado.

Los archivos automáticos \texttt{latencias\_*.csv} y \texttt{eventos\_*.csv} se mantuvieron separados de las mediciones consolidadas, ya que contienen también registros de calentamiento y pruebas funcionales.

\subsection{Limitaciones y amenazas a la validez}

Las mediciones se realizaron en localhost, por lo que no incorporan las condiciones de una comunicación entre computadoras a través de una red física. Factores como la latencia de red, la congestión o la variabilidad de los enlaces podrían modificar los resultados en otro entorno.

Aunque ambos mecanismos utilizaron las mismas cargas útiles, sus operaciones internas no fueron idénticas. En TCP, la serialización JSON se realizó dentro del intervalo medido. En gRPC, el objeto de solicitud se construyó antes de iniciar el cronómetro y la llamada remota quedó incluida en la medición. En consecuencia, las diferencias observadas corresponden a las implementaciones evaluadas y no pueden atribuirse exclusivamente al protocolo de transporte.

También deben distinguirse las métricas utilizadas. El RTT individual no incluye todas las operaciones de registro posteriores a la recepción del acuse, mientras que el rendimiento considera la duración completa de cada corrida, incluyendo dichas operaciones.

Las pruebas TCP se ejecutaron antes que las de gRPC. No se controló instrumentalmente la actividad de otros procesos de Windows ni la variación de carga del procesador. La planificación de tareas del sistema operativo y otras actividades en segundo plano pudieron influir en los tiempos obtenidos.

Por estas razones, los resultados permiten comparar el comportamiento de los clientes implementados por el equipo bajo las condiciones específicas de esta práctica, sin establecer conclusiones generales sobre el rendimiento de TCP y gRPC en cualquier sistema distribuido.